% ════════════════════════════════════════════════════════════════
%  CHAPITRE 1 — CADRE GÉNÉRAL DU PROJET
% ════════════════════════════════════════════════════════════════

\chapter{Cadre général du projet}

% ── Introduction du chapitre ────────────────────────────────────
\section*{Introduction}
\addcontentsline{toc}{section}{Introduction}

Ce chapitre fixe le cadre du projet : organisme d'accueil, situation
initiale, périmètre et méthodologie Scrum retenue pour conduire la
réalisation.

% ════════════════════════════════════════════════════════════════
\section{Présentation de l'organisme d'accueil}
% ════════════════════════════════════════════════════════════════

\subsection{Historique}

\textbf{Schulte \& Co. GmbH} est un groupe industriel allemand fondé à
Lüdenscheid. L'entreprise s'est développée dans le secteur des composants
électriques et électroniques pour l'automobile.

Pour accompagner la croissance de ses activités, le groupe a ouvert une
filiale en Tunisie : \textbf{Schulte Automotive Tunisia Sarl}. Cette
implantation répond à un objectif clair : rapprocher la production des
marchés européens tout en s'appuyant sur des compétences locales.

\subsection{Localisation}

La filiale tunisienne opère sur deux sites :
\begin{itemize}
    \item \textbf{Bouarada} (gouvernorat de Siliana) ;
    \item \textbf{Zaghouan}.
\end{itemize}

Ce fonctionnement multi-sites influence directement l'organisation RH,
car chaque site gère ses besoins de recrutement selon son activité et
son planning de production.

\subsection{Domaines d'activité}

Le groupe Schulte intervient principalement dans la fabrication de
composants destinés à l'industrie automobile, notamment :
\begin{itemize}
    \item l'assemblage de câbles ;
    \item les faisceaux et systèmes de connexion ;
    \item des sous-ensembles électriques pour constructeurs automobiles.
\end{itemize}

L'environnement industriel impose un niveau élevé de qualité, de
traçabilité et de réactivité, y compris dans les processus RH.

\subsection{Département d'accueil}

Notre stage s'est déroulé au sein du \textbf{département Ressources
Humaines} de Schulte Automotive Tunisia.

Ce département gère, entre autres, la publication des offres d'emploi,
la réception des candidatures, la présélection des profils et la
planification des entretiens. C'est donc le service le plus concerné
par une démarche de digitalisation du recrutement.

Cette proximité avec les opérations quotidiennes RH a permis d'ancrer le
projet dans des besoins concrets, observés et validés sur le terrain.

% ════════════════════════════════════════════════════════════════
\section{Présentation du projet}
% ════════════════════════════════════════════════════════════════

\subsection{Contexte}

Avant ce projet, la gestion des candidatures reposait encore sur des
pratiques manuelles. Plusieurs candidats déposaient leur CV en papier
ou l'envoyaient par courrier, sans canal numérique local dédié.

En parallèle, le site web du groupe existe, mais il ne répond pas aux
besoins réels du recrutement en Tunisie. Les offres visibles concernent
majoritairement d'autres pays et le parcours de candidature n'est pas
adapté au contexte local.

\subsection{Problématique}

La problématique centrale est la suivante :
\textit{comment mettre en place un processus de recrutement numérique,
clair et traçable, adapté aux deux sites de Schulte Tunisia ?}

Cette problématique se traduit par trois besoins immédiats :
\begin{itemize}
    \item améliorer l'accès aux offres pour les candidats ;
    \item structurer le suivi des candidatures côté RH ;
    \item assurer une communication plus claire sur l'état des dossiers.
\end{itemize}

\subsection{Périmètre}

Le périmètre de notre travail couvre la conception et la réalisation
d'une plateforme de recrutement dédiée à Schulte Automotive Tunisia.

La solution comprend :
\begin{itemize}
    \item une interface \textbf{candidat} (PWA) pour consulter les offres,
    postuler et suivre l'évolution de la candidature ;
    \item une interface \textbf{RH} pour publier les offres, analyser les
    profils et gérer le pipeline de recrutement ;
    \item une interface \textbf{administration} pour gérer les comptes RH
    et les modèles d'offres.
\end{itemize}

Le projet cible les deux sites tunisiens (Bouarada et Zaghouan) avec une
séparation des données entre les équipes RH.

\subsection{Analyse de l'existant}

L'analyse de la situation initiale a montré un fonctionnement peu outillé :
\begin{itemize}
    \item réception des candidatures par papier, email ou courrier ;
    \item tri manuel des CV ;
    \item absence de tableau de suivi partagé ;
    \item planification des entretiens sans workflow unifié ;
    \item retour irrégulier vers les candidats.
\end{itemize}

Ce mode de travail peut fonctionner avec un faible volume, mais devient
vite difficile à maintenir quand le nombre de candidatures augmente.

\subsection{Critique de l'existant}

À partir de cette analyse, nous retenons les limites suivantes :
\begin{itemize}
    \item manque de transparence pour les candidats ;
    \item charge administrative élevée pour le service RH ;
    \item risque d'erreurs et de perte d'information ;
    \item absence de traçabilité complète des décisions ;
    \item difficulté de pilotage entre les deux sites.
\end{itemize}

Ces limites justifient la mise en place d'une solution numérique adaptée
au contexte local de l'entreprise, avec une gouvernance claire des accès
et une meilleure continuité du suivi candidat.

\subsection{Solution proposée}

La solution proposée est une plateforme web de recrutement, pensée pour
les besoins de Schulte Tunisia.

Elle permet de :
\begin{itemize}
    \item centraliser toutes les offres et candidatures ;
    \item suivre chaque candidature par étapes (nouveau, en examen,
    entretien, accepté, refusé) ;
    \item notifier les acteurs concernés au bon moment ;
    \item améliorer la qualité du traitement RH grâce à des outils de
    tri, de suivi et d'évaluation.
\end{itemize}

Cette solution vise un recrutement plus rapide, plus lisible et plus
équitable pour les candidats comme pour l'entreprise.

Elle ne se limite pas à des formulaires de saisie : elle organise un
processus multi-acteurs avec règles de sécurité, gestion d'états et
notifications, afin de fiabiliser le cycle de recrutement de bout en bout.

La section suivante explique la méthode de travail utilisée pour passer
de ce besoin métier à une implémentation concrète et progressive.

% ════════════════════════════════════════════════════════════════
\section{Méthodologie de travail}
% ════════════════════════════════════════════════════════════════

\subsection{Méthodes agiles}

Nous avons adopté une approche agile, car elle permet de livrer le projet
par étapes courtes et d'ajuster les priorités selon les retours terrain.

Dans ce cadre, chaque itération produit un résultat concret : une
fonctionnalité testable, puis consolidée dans l'itération suivante.
Cette logique réduit le risque de décalage entre la solution développée
et le besoin métier réel.

\subsection{Scrum (rôles, artefacts, événements)}

\textbf{Rôles :}
\begin{itemize}
    \item \textbf{Product Owner} : rôle assuré principalement par
    l'encadrant et les responsables RH, qui expriment les besoins et
    valident les priorités ;
    \item \textbf{Scrum Master} : rôle assuré par l'étudiant pour
    organiser le travail, lever les blocages et garder le rythme ;
    \item \textbf{Équipe de développement} : dans notre cas,
    développement réalisé en mode solo.
\end{itemize}

\textbf{Artefacts :}
\begin{itemize}
    \item \textbf{Product Backlog} : liste priorisée des fonctionnalités ;
    \item \textbf{Sprint Backlog} : tâches retenues pour chaque sprint ;
    \item \textbf{Incrément} : version fonctionnelle livrée à la fin de
    chaque sprint.
\end{itemize}

\textbf{Événements :}
\begin{itemize}
    \item \textbf{Sprint Planning} : définition des objectifs du sprint ;
    \item \textbf{Daily Scrum} : suivi quotidien de l'avancement
    (adapté à une organisation individuelle) ;
    \item \textbf{Sprint Review} : démonstration des éléments livrés et
    collecte des retours ;
    \item \textbf{Sprint Retrospective} : identification des points à
    améliorer pour le sprint suivant.
\end{itemize}

% ── Conclusion du chapitre ──────────────────────────────────────
\section*{Conclusion du chapitre 1}
\addcontentsline{toc}{section}{Conclusion du chapitre 1}

Ce chapitre a permis de cadrer précisément le projet : environnement
industriel, besoins RH, limites de l'existant et périmètre de la solution.
Il justifie également le choix d'une démarche Scrum, cohérente avec un
développement progressif et orienté valeur.

Sur cette base, le chapitre suivant présente l'état de l'art, puis les
choix technologiques retenus pour concrétiser la solution.
